\section{From Mathematics to Code: Three-Tier Derivation}\label{sec:m2c}

\subsection{Tier 1: The Perturbation Operator (Mathematics)}

\cref{def:perturbation} is the mathematical contract. The five
invariants F1--F5 (\cref{ssec:invariants}) are the obligations the
next tiers must preserve.

\subsection{Tier 2: The Algorithm}

The perturbation operator translates to a deterministic
kernel-conformant algorithm. A \texttt{pre\_handler} is registered on
each kprobe (one for \texttt{submit\_bh}, one for
\texttt{submit\_bio\_noacct}); on each invocation, the handler reads
the relevant arguments, applies the filter, samples the PRNG, and
performs the flip via \texttt{kmap\_local\_page} for arm64-safe page
access~\cite{kernel-highmem}.

\begin{lstlisting}[caption={RadFI bio pre-handler kprobe on
\texttt{submit\_bio\_noacct}.},label={lst:prehandler}]
/*
 * RadFI bio pre-handler -- kprobe on submit_bio_noacct.
 * Returns 0 always (pre-handler does not redirect execution).
 */
static int radfi_kp_submit_bio_pre(struct kprobe *kp,
                                   struct pt_regs *regs)
{
    struct radfi_state *st = radfi_global;
    struct bio *bio;
    unsigned int op;
    dev_t dev;

    if (!st || !READ_ONCE(st->hook_blk))
        return 0;

    bio = (struct bio *)regs->regs[0];
    if (!bio || !bio->bi_bdev)
        return 0;

    /* F3: filter monotonicity (op type) */
    op = bio_op(bio);
    if (op == REQ_OP_WRITE) {
        /* always candidate */
    } else if (op == REQ_OP_READ) {
        if (!READ_ONCE(st->inject_on_read))
            return 0;        /* v0.1.2 flag */
    } else {
        return 0;            /* skip non-data ops */
    }

    /* F3: filter monotonicity (device) */
    dev = bio->bi_bdev->bd_dev;
    if (READ_ONCE(st->target_dev) &&
        (u32)dev != READ_ONCE(st->target_dev))
        return 0;

    atomic64_inc(&st->call_count);

    /* F1: probability sampling */
    if (!radfi_should_inject(st)) {
        atomic64_inc(&st->skipped_prob);
        return 0;
    }

    /* F1, F2: single bit-flip on first segment, uniform */
    radfi_flip_one_bit_in_first_segment(bio);

    atomic64_inc(&st->flip_count);
    return 0;
}
\end{lstlisting}

The listing above shows the principal control flow; the actual
implementation includes additional safety checks on
\texttt{bio->bi\_vcnt} and \texttt{bio->bi\_iter.bi\_size},
omitted here for clarity.

The algorithm is structurally simple. Its correctness rests entirely
on the alignment between the mathematical invariants and the
kernel-side implementation choices. \cref{tab:m2c-trace} makes this
alignment explicit.

\subsection{Tier 3: The Kernel Implementation}

The reference implementation is a Linux kernel module
(\texttt{radfi.ko}, GPL-2.0) registering two kprobes: one on
\texttt{submit\_bh} (filesystem hook), one on
\texttt{submit\_bio\_noacct} (block layer hook). Configuration is
exposed through debugfs entries under
\texttt{/sys/kernel/debug/radfi/}:

\begin{itemize}
\item \texttt{enabled} (bool): master switch.
\item \texttt{hook\_fs} (bool), \texttt{hook\_blk} (bool): per-hook
  enable.
\item \texttt{inject\_on\_read} (bool, v0.1.2): DRAM-bidirectional
  toggle.
\item \texttt{probability} (u32): PPM probability.
\item \texttt{target\_dev} (u32), \texttt{target\_inode} (u64),
  \texttt{target\_block} (u64): filter values.
\item \texttt{seed} (u64): PRNG seed.
\item \texttt{call\_count}, \texttt{flip\_count},
  \texttt{skipped\_disabled}, \texttt{skipped\_prob} (read-only u64):
  F5 counters.
\end{itemize}

The kmap path uses \texttt{kmap\_local\_page} (rather than the
deprecated \texttt{kmap\_atomic}) for arm64
safety~\cite{kernel-highmem}. The PRNG state is per-module-load,
re-seeded at module init. In the actual implementation, fields
accessed multiple times in the pre-handler (e.g.\ \texttt{target\_dev})
are loaded once into a local variable to avoid double
\texttt{READ\_ONCE} hazards under concurrent debugfs writes; the
listing elides this for readability.

\paragraph{Differences from generic SDC fault-injection tools.}
Tools such as Xception~\cite{carreira1998xception} and
EDFI~\cite{giuffrida2013edfi} inject into instruction streams or
memory regions selected by static analysis. RadFI restricts itself
to the bio-layer payload, which corresponds precisely to the
abstraction at which a filesystem can defend (the BEAMFS recovery
operator \(\mathcal{G}\) acts on bytes read back from disk, not on
instructions). This narrow scope is what makes the
falsification-of-recovery-theorems use case crisp.

\paragraph{Choice of kprobes over fprobe.}
The v0.1.x build uses \texttt{register\_kprobe}\cite{kernel-kprobes}
for compatibility with Linux kernel 7.0.x and earlier; migration
to the newer \texttt{fprobe} interface is v2 work.

\paragraph{Architecture-specific register access.}
The pre-handler reads its bio argument from \texttt{regs->regs[0]},
which is the aarch64 first-argument register \texttt{x0}. The
x86\_64 path uses \texttt{regs->di}; the production module selects
the correct field via \texttt{\#ifdef CONFIG\_ARM64} /
\texttt{CONFIG\_X86\_64}.

\paragraph{Allocation contexts.}
The kprobe pre-handler runs in interrupt context (the bio submission
path may be invoked under spinlocks or in softirq); all allocations
on the pre-handler path use \texttt{GFP\_ATOMIC}. The slow paths
(module init, debugfs writes) use \texttt{GFP\_KERNEL}.

\paragraph{Kprobe blacklist verification.}
Neither \texttt{submit\_bio\_noacct} nor \texttt{submit\_bh} appears
in the \texttt{\_\_kprobes} blacklist as of Linux kernel 7.0.x; this
is verified at module load by inspecting the kprobe
registration return value.

\paragraph{Singleton state and replay determinism.}
RadFI uses a singleton state (\texttt{radfi\_global}) rather than a
per-CPU structure. This is intentional: a per-CPU PRNG state would
break F4 (replay determinism) by making the flip sequence depend on
which CPU happened to service each bio.

\paragraph{Module unload safety.}
Module unload calls \texttt{unregister\_kprobe} followed by
\texttt{synchronize\_rcu}, ensuring any in-flight pre-handler
invocation completes before the module's data structures are freed.

\paragraph{Production safety.}
The v0.1.x build aborts module load if the host system has any
non-loop block device mounted with a BEAMFS filesystem. This is a build-time check enforced by a runtime probe
in \texttt{radfi\_init}. The intent is to make accidental injection
on a production filesystem impossible without explicit code
modification.

\subsection{Traceability: Math \(\to\) Algorithm \(\to\) Code}

\cref{tab:m2c-trace} expresses the descent in tabular form, in the
same auditable contract style
as~\cite[Table~II]{desbrieres2026beamfsv1}.

\begin{table*}[t]
\caption{Traceability of the perturbation operator. Each
mathematical commitment maps to an algorithmic step and a kernel
construct. The ``Phase'' column indicates when the construct is
checked: Mount = once at module load; Runtime = on each bio
interception.}
\label{tab:m2c-trace}
\centering
\small
\begin{tabular}{@{}p{3.2cm}p{4.6cm}p{6.5cm}l@{}}
\toprule
Math Commitment (Tier 1) & Algorithm (Tier 2) & Kernel Construct (Tier 3) & Phase \\
\midrule
F1 (Single-bit primitive) & \texttt{radfi\_flip\_one\_\allowbreak bit\_in\_\allowbreak first\_segment}\\flips one bit per call & Sampling: \texttt{byte\_off = prng \% bv.bv\_len}; \texttt{bit\_pos = prng \% 8}; XOR & Runtime \\
F2 (Targeting independence) & Sample positions from bio geometry only, not from byte values & \texttt{bv.bv\_len}, \texttt{bv.bv\_offset} used as bounds; payload bytes never read & Runtime \\
F3 (Filter monotonicity) & Sequential filter checks return 0 on any non-match & \texttt{if ((u32)dev != target\_dev) return 0;} per filter component & Runtime \\
F4 (Replay determinism) & xoshiro256** seeded by \texttt{seed} parameter & \texttt{radfi\_prng\_state} initialized from \texttt{st->prng\_seed} at module init & Mount \\
F5 (Counter accounting) & \texttt{atomic64\_inc} on each path & \texttt{\&st->call\_count}, \texttt{\&st->flip\_count}, \texttt{\&st->skipped\_*} & Runtime \\
\cref{def:perturbation}, step 1 (Intercept) & kprobe \texttt{pre\_handler} & \texttt{register\_\allowbreak kprobe(\&radfi\_kp\_\allowbreak submit\_bio)} & Mount \\
\cref{def:perturbation}, step 2 (Filter) & Filter algebra on bio attributes & \texttt{bio\_op(bio)}, \texttt{bio->bi\_bdev->bd\_dev} & Runtime \\
\cref{def:perturbation}, step 3 (Probability) & PRNG comparison vs \(\mathit{prob}\) & \texttt{radfi\_should\_inject(st)} & Runtime \\
\cref{def:perturbation}, step 4 (Flip) & \texttt{kmap\_local\_page} + XOR + \texttt{kunmap\_local} & arm64-safe page access~\cite{kernel-highmem} & Runtime \\
\cref{def:filter} (Filter algebra) & Tuple of optional filters & debugfs entries \texttt{target\_dev}, \texttt{target\_inode}, \texttt{target\_block}, op type & Mount \\
DRAM-bidirectional (\cref{ssec:dram-bidir}) & \texttt{inject\_on\_read} toggle on \texttt{REQ\_OP\_READ} & \texttt{READ\_ONCE(st->inject\_on\_read)} at op type filter & Runtime \\
\bottomrule
\end{tabular}
\end{table*}
